\SourceLeaf{252}%
\(\mathrm{d}p,\mathrm{d}q,\mathrm{d}p',\mathrm{d}q'\) entsprechen,
\[
\BareRoot{E}\mathord{.}\mathrm{d}p,\qquad\BareRoot{G}\mathord{.}\mathrm{d}q,\qquad\BareRoot{E'}\mathord{.}\mathrm{d}q',\qquad\BareRoot{G'}\mathord{.}\mathrm{d}q'
\]
Die Winkel, welche die Richtungen dieser Elemente mit einer beliebigen festen Richtung bilden, bezeichnen wir mit \(M,N,M',N'\). Wir messen sie in dem Sinn, in dem die zweite Richtung zur ersten liegt, sodass \(\sin(N-M)\) positiv ist. Wir nehmen an (was zulässig ist), dass die vierte im selben Sinn zur dritten liegt, sodass auch \(\sin(N'-M')\) positiv ist. Betrachten wir unter diesen Festlegungen einen anderen Punkt, der vom ersten unendlich wenig entfernt ist und dem die Werte der Veränderlichen
\[
p+\mathrm{d}p,\quad q+\mathrm{d}q,\quad p'+\mathrm{d}p',\quad q'+\mathrm{d}q'
\]
entsprechen, so erkennt man bei ein wenig Aufmerksamkeit, dass allgemein, also unabhängig von den Werten der Änderungen \(\mathrm{d}p,\mathrm{d}q,\mathrm{d}p',\mathrm{d}q'\), gilt
\[
\begin{aligned}
&\BareRoot{E}\mathord{.}\mathrm{d}p\mathord{.}\sin M+\BareRoot{G}\mathord{.}\mathrm{d}q\mathord{.}\sin N\\
&=\BareRoot{E'}\mathord{.}\mathrm{d}p'\mathord{.}\sin M'+\BareRoot{G'}\mathord{.}\mathrm{d}q'\mathord{.}\sin N'
\end{aligned}
\]
denn jeder der beiden Ausdrücke ist einfach der Abstand des neuen Punkts von der Geraden, von der aus die Richtungswinkel gemessen werden. Nach der oben eingeführten Schreibweise gilt aber \(N-M=\omega\). Entsprechend setzen wir \(N'-M'=\omega'\) und außerdem \(N-M'=\psi\). Die soeben gefundene Gleichung lässt sich daher in der folgenden Form schreiben
\[
\begin{aligned}
&\BareRoot{E}\mathord{.}\mathrm{d}p\mathord{.}\sin(M'-\omega+\psi)+\BareRoot{G}\mathord{.}\mathrm{d}q\mathord{.}\sin(M'+\psi)\\
&=\BareRoot{E'}\mathord{.}\mathrm{d}p'\mathord{.}\sin M'+\BareRoot{G'}\mathord{.}\mathrm{d}q'\mathord{.}\sin(M'+\omega')
\end{aligned}
\]
oder in dieser Form
\[
\begin{aligned}
&\BareRoot{E}\mathord{.}\mathrm{d}p\mathord{.}\sin(N'-\omega-\omega'+\psi)+\BareRoot{G}\mathord{.}\mathrm{d}q\mathord{.}\sin(N'-\omega'+\psi)\\
&=\BareRoot{E'}\mathord{.}\mathrm{d}p'\mathord{.}\sin(N'-\omega')+\BareRoot{G'}\mathord{.}\mathrm{d}q'\mathord{.}\sin N'
\end{aligned}
\]

Da die Gleichung offenbar von der Ausgangsrichtung unabhängig sein muss, darf man diese beliebig wählen. Setzen wir in der zweiten Form \(N'=0\) oder in der ersten \(M'=0\), so erhalten wir die folgenden Gleichungen:
\[
\begin{aligned}
\BareRoot{E'}\mathord{.}\sin\omega'\mathord{.}\mathrm{d}p'&=\BareRoot{E}\mathord{.}\sin(\omega+\omega'-\psi)\mathord{.}\mathrm{d}p+\BareRoot{G}\mathord{.}\sin(\omega'-\psi)\mathord{.}\mathrm{d}q\\
\BareRoot{G'}\mathord{.}\sin\omega'\mathord{.}\mathrm{d}q'&=\BareRoot{E}\mathord{.}\sin(\psi-\omega)\mathord{.}\mathrm{d}p+\BareRoot{G}\mathord{.}\sin\psi\mathord{.}\mathrm{d}q
\end{aligned}
\]
Da diese Gleichungen identisch sein müssen mit
\[
\begin{aligned}
\mathrm{d}p'&=\alpha\mathrm{d}p+\beta\mathrm{d}q\\
\mathrm{d}q'&=\gamma\mathrm{d}p+\delta\mathrm{d}q
\end{aligned}
\]
